\chapter{文献分层与继承边界}
\label{ch:lit}

模型中的每一层都标明来源。继承表示后文直接使用其方程结构；改写表示边界、状态或判据换成了热加工决策；不继承表示该文献的结论不能被写成本文的结果。条目只收录可以核到作者、题目和年份的来源。正式引文的卷期页码在成稿参考文献中逐条核对，不在大纲里用印象填写。

\section{多孔介质与食品热湿传递}

\begin{table}[H]
\centering
\caption{传递骨架}
\begin{tabularx}{\textwidth}{lXll}
\toprule
来源 & 使用的结果 & 方式 & 不继承 \\
\midrule
Whitaker, 1977 & 体积平均后的能量、动量与组分方程 & 继承 & 干燥工艺的具体边界 \\
Datta, 2007 & 把上述骨架用于食品同时热质传递 & 继承 & 某一烘箱的实验拟合 \\
Ni、Datta 等关于挤压与油炸多孔介质的系列 & 蒸发—冷凝源项、压力驱动汽流 & 改写到蒸烤界面 & 油炸外部流动 \\
Feyissa、Gernaey、Adler-Nissen, 2013 & 对流烤箱中肉块的三维耦合场、冷点可偏离几何中心 & 继承现象，方程改为一维主算例加二维对照 & 该文的网格与物料参数数值 \\
\bottomrule
\end{tabularx}
\end{table}

由此得到 \eqref{eq:energy}--\eqref{eq:phase} 的结构。蒸比干热更快进入内部，写成冷凝潜热，不写成“蒸汽导热更快”。干壳写成近表面渗透率与 $k_{\mathrm{eff}}$ 下降。

\section{水分活度、收缩与各向异性}

\begin{table}[H]
\centering
\caption{界面与结构闭合}
\begin{tabularx}{\textwidth}{lXll}
\toprule
来源 & 使用的结果 & 方式 & 不继承 \\
\midrule
van den Berg、Bruin 关于食品水分等温线的综述传统；GAB 形式 & $a_w$ 与含水率的非线性关系 & 继承形式 \eqref{eq:gab} & 未经测定的 $M_0,C,K$ 数值 \\
Clausius--Clapeyron & 吸附热使等温线随温度移动 & 继承 & 把活度当成含水率本身 \\
肉类热导率综述（如 Choi、Okos 的组分加和传统） & $k$ 随水分、脂肪、温度变化 & 继承趋势 & 把各向同性标量当成全部 \\
纤维食品与肌肉各向异性测量文献 & $k_\parallel\neq k_\perp$ & 写成 \eqref{eq:aniso} & 某一品种的单点数据当普遍常数 \\
蛋白受热收缩与孔隙变化的工艺文献 & $L$ 与 $K$ 随 $\alpha_p$ 变化 & 写成 \eqref{eq:shrink} & 收缩即品质的定义 \\
\bottomrule
\end{tabularx}
\end{table}

$M_0,C,K,k_\parallel/k_\perp,\beta_s$ 在大纲中只作为有区间的参数。成稿里每个数旁边要有文献区间；没有区间的数不得进入主结果，只能进入灵敏度扫描。

\section{反应进度与热力当量}

\begin{table}[H]
\centering
\caption{进度，而不是温度阈值}
\begin{tabularx}{\textwidth}{lXll}
\toprule
来源 & 使用的结果 & 方式 & 不继承 \\
\midrule
van Boekel, 2008 & 食品品质动力学用速率方程，Arrhenius 只是其中一种温度依赖 & 继承速率方程；温度依赖仍用 Arrhenius，并在讨论中保留其适用范围 & 把所有品质都写成一级反应 \\
Holdsworth 关于热力杀菌的标准表述 & $F$ 值沿冷点历史积分 & 写成 \eqref{eq:F} 的 $\alpha_m$ & 把 $\alpha_{\mathrm{safe}}$ 写成法规限值 \\
蛋白变性、淀粉糊化、美拉德的动力学实验文献 & 三类进度对水分的开关不同 & $g_i(\varepsilon_l,S)$ & 用颜色代替中心进度 \\
\bottomrule
\end{tabularx}
\end{table}

美拉德文献解释表面色泽路径 B。它不提供冷点是否安全的证据。微生物文献只提供积分形式，初始污染量不在状态向量中。

\section{逆热传导与可观测性}

\begin{table}[H]
\centering
\caption{为什么表面温度不够}
\begin{tabularx}{\textwidth}{lXll}
\toprule
来源 & 使用的结果 & 方式 & 不继承 \\
\midrule
Özisik、Orlande, 2000 & 逆热传导对高频噪声指数放大 & 继承病态结构 & 他们的正则化算法作为唯一估计器 \\
Beck 关于序贯估计的传统 & 未来时间步可稳定热流估计 & 作为病态程度的参照，不作为在线主算法 & 把整机问题降成单层壁面 \\
线性系统可观测性与 Kalman 滤波 & 秩条件、协方差更新 & 用于两节点命题 & 线性模型覆盖相变 \\
Evensen, 2003 & 集合卡尔曼滤波的预测—分析格式 & 场层估计 & 把集合均值当成决策量 \\
\bottomrule
\end{tabularx}
\end{table}

命题 \ref{prop:steady} 与 \ref{prop:pulse} 是本文在两节点蒸烤边界上的推论，不是上述文献的原定理。成稿中这两条命题的证明自洽，文献只为病态性和集合更新负责。

\section{信息价值与对偶控制}

\begin{table}[H]
\centering
\caption{动作为什么可以是一次实验}
\begin{tabularx}{\textwidth}{lXll}
\toprule
来源 & 使用的结果 & 方式 & 不继承 \\
\midrule
Lindley, 1956 & 实验的信息量是先验与后验熵之差 & $G(a)$ 的定义 & 把效用只写成参数熵 \\
Fedotov、Kiefer 等最优实验设计 & $\log\det$ 与参数置信椭球 & Fisher 形式的对照 & 离线一次性设计代替在线脉冲 \\
Feldbaum 对偶控制 & 控制量同时影响状态与信息 & \eqref{eq:dual} 的问题提法 & 其原问题的解析解 \\
Tse、Bar-Shalom 等谨慎控制与主动学习控制 & 有限时域上信息与代价并列 & 探针库是该问题的可实现子集 & 把家电执行器当成任意波形 \\
\bottomrule
\end{tabularx}
\end{table}

创新边界在这里：信息目标不是新的熵定义，而是与相变变号、对流切换绑定的可执行子集，以及“信息未下降则停止时刻不动”的判定。

\section{算子学习}

\begin{table}[H]
\centering
\caption{加速器，不是定律}
\begin{tabularx}{\textwidth}{lXll}
\toprule
来源 & 使用的结果 & 方式 & 不继承 \\
\midrule
Raissi、Perdikaris、Karniadakis, 2019 & 残差可惩罚方程 & 训练时的物理损失项 & 用 PINN 代替在线安全判断 \\
Li 等, ICLR 2021 & Fourier 神经算子学习参数化演化 & 候选 $\mathcal{G}_\phi$ & 训练分布之外的精度保证 \\
Lu 等, 2021 & DeepONet 的分支—主干结构 & 同一数据上的对照算子 & 一种结构优于另一种的普遍结论 \\
\bottomrule
\end{tabularx}
\end{table}

主模型的选择规则已经写明：同一训练集上，蒸汽干预位移误差更低者为主模型。文献不提供这个选择的胜负。

\section{因果、预测控制与安全过滤}

\begin{table}[H]
\centering
\caption{解释与执行}
\begin{tabularx}{\textwidth}{lXll}
\toprule
来源 & 使用的结果 & 方式 & 不继承 \\
\midrule
Pearl, 2009 & 预测、干预、反事实的分层；do 是替换结构方程 & 路径 A、B、C 的计算顺序 & 从短烹饪序列里发现父节点 \\
Rawlings、Mayne、Diehl, 2017 & 有限时域约束优化 & \eqref{eq:mpc} 的问题类 & 某一求解器的默认最优性声明 \\
Ames、Xu、Grizzle、Tabuada, 2017 & 屏障函数约束下的二次投影 & \eqref{eq:cbf} 的形式 & 屏障条件自动等于器具安全认证 \\
IEC 60335-1 与 GB 4706.1 & 漏电、温升、异常工作的器具安全框架 & 硬回路的要求来源 & 把标准条款写成算法已满足 \\
\bottomrule
\end{tabularx}
\end{table}

命题 \ref{prop:shield} 使用的是“可行域自变量与后验不交”。它是投影形式加上测量隔离之后的直接推论。标准负责硬回路必须存在，不负责软件投影的参数。

\section{创新所在的层}

把上面各表的“不继承”列合起来，本文新增的不是多孔介质方程、熵、do 算子、神经算子或屏障函数。新增的是四条连接：

\begin{enumerate}
\item 决策状态 $\vect{z}$ 从全场中取出，并证明稳态腔温不含 $R_{\mathrm{int}}$。
\item 信息动作落成相变变号后接对流切换，且无效则不移动停止时刻。
\item 同一蒸汽干预分解为潜热、表面品质、结壳热阻三条路径。
\item 品质后验与安全可行域使用不相交的测量。
\end{enumerate}

这四条若与某篇已发表论文的技术方案相同，成稿时改为引用并退出创新表述。大纲阶段不把“未读到”写成“不存在”。
